\section{The Framework}\label{sec:framework}

The formalisation is built on an Agda library for presentations of
monoids and groups by generators and relations, normal forms, and
completeness of finite relation sets, with wire-indexed circuits as its
motivating instance~\cite{qupit-code,bian2026framework}.  This section recalls, briefly and
with code quoted verbatim, the parts the rest of the paper uses.  The
emphasis is on two things: the generic lemma
\aid{by-normalization}, which is the only place completeness is ever
proved, and the \emph{presentation construction theorems} --- for
semidirect products and for central extensions --- which are what allow
\S\S\ref{sec:symplectic}--\ref{sec:composing} to deal with
presentations and normal forms of the symplectic group, the Pauli group
and the scalars one at a time.

\subsection{Words and presented monoids}\label{sec:presentations}

A word over a generator set $X$ is an unquotiented syntax tree with
constructors $[\_]\aid{ʷ}$, \aid{ε} and \aid{•} --- deliberately not a
list, so that a word remembers its bracketing and rules may be applied
to any parenthesised subterm.  A relation set is $\aid{WRel}\;X = \aid{Rel}\;(\aid{Word}\;X)\;0\ell$.
Given $\Gamma : \aid{WRel}\;X$, the library's two-level convention is
that \aid{{\_}==={\_}} always means the raw axioms and \aid{{\_}≈{\_}}
the monoid congruence they generate, an inductive family closed under
reflexivity, symmetry, transitivity, congruence of \aid{•},
associativity, the unit laws, and \aid{axiom}.  Thus $\approx$ \emph{is}
the word problem of $\langle X \mid \Gamma\rangle$, and a derivation is
a first-class object one can induct on.  Because words have no primitive
inverses every presentation is at bottom a monoid presentation; a
\aid{Grouplike} witness (every generator has a left inverse up to
$\approx$) upgrades the presented monoid to the presented group
\aid{•-ε-group}.  The extension of a generator map $f : X \to
\aid{Word}\;Y$ to words is written postfix, $(f\,\aid{ʷ})$, and a coset
action $h : C \to Y \to \aid{Word}\;X \times C$ extends to words as
$(h\,\aid{ᵗ})$, threading the coset from left to right; these two
combinators, with the conjugation extensions $(\mathit{conj}\,\aid{ʰ'})$
and $(\mathit{conj}\,\aid{ⁿ'})$, are the whole syntactic vocabulary of
the framework.

A presentation of a concrete group is a record:

\begin{agdacode}
record _IsPresentationOf_ (_===_ : WRel X) (G : Group a ℓ) : Set (a ⊔ ℓ) where
  field
    gl : Grouplike _===_
  module GL = Group-Lemmas _===_ gl
  open GroupMorphisms (Group.rawGroup GL.•-ε-group) (Group.rawGroup G)
  field
    ⟦_⟧ : Word X → Group.Carrier G
    iso : IsGroupIsomorphism ⟦_⟧
\end{agdacode}

A weaker record, \aid{{\_}IsSubPresentationOf{\_}}, asks only for a
monomorphism; it packages soundness (the homomorphism's congruence) and
completeness (its injectivity), both stated with the standard library's
\aid{Congruent} and \aid{Injective}, and \aid{isPresentationOf} upgrades
it to a presentation given surjectivity.  Every presentation theorem in
this paper is assembled this way: soundness, completeness by
normalisation, and surjectivity are proved separately and the records do
the rest.

\subsection{Normal forms, and completeness by normalisation}\label{sec:nfproperty}

Normal-form witnesses are, on the nose, the standard library's function
bundles out of the word setoid: $\aid{NormalFormInjective}\;\Gamma\;\NF$
is an \aid{Injection}, $\aid{NormalForm}\;\Gamma\;\NF$ a
\aid{RightInverse} --- the map \aid{nf} together with a section
$\aid{inv-nf} : \NF \to \aid{Word}\;X$ and a proof $\aid{inv-nf} \circ
\aid{nf} \approx \mathrm{id}$ --- and $\aid{BijectiveNormalForm}$ a
\aid{Bijection}.  The section is data: it is the algorithm that realises
each normal form as a canonical word, and it is what completeness
consumes.  Fix a semantics $\sem{\_} : \aid{Word}\;X \to \mathit{Sem}$
into a setoid.  A normal form is \emph{unique for} $\sem{\_}$ when
representatives with equal denotations are equal normal forms, and
\emph{soundness plus a unique normal form yields completeness}:

\begin{agdacode}
record UniqueNormalForm (inv-nf : |NF| → Word X) : Set (c ⊔ d) where
  field
    unique : ∀ {u v : |NF|} → ⟦ inv-nf u ⟧ ≈₂ ⟦ inv-nf v ⟧ → u ≈ₙ v

by-normalization : {normalForm : NormalForm} → let open NormalForm normalForm in
  UniqueNormalForm inv-nf → Congruent _≈_ _≈₂_ ⟦_⟧ → Injective _≈_ _≈₂_ ⟦_⟧
\end{agdacode}

The proof is five lines: given $\sem{w} \approx_2 \sem{v}$, soundness
transports along $w \approx \aid{inv-nf}\,(\aid{nf}\,w)$, uniqueness
identifies the normal forms, and injectivity of \aid{nf} closes the
loop.  Its value is that every case study only ever proves
\aid{unique} --- a statement about normal forms, i.e.\ about concrete
first-order data --- never about arbitrary words.  This is precisely the
division of labour of the qupit paper's Proposition~3.8 (uniqueness of
the symplectic normal form) and Proposition~4.3 (every circuit rewrites
to normal form), except that the rewriting half is packaged as a
\emph{function} with a verified section, and the argument
``two circuits with the same symplectic matrix rewrite to the same
normal form, hence to each other'' is discharged once and for all.  A
converse, \aid{by-completeness}, recovers \aid{UniqueNormalForm} from
completeness plus an exact section $\aid{nf} \circ \aid{inv-nf} \equiv
\mathrm{id}$; it is the route by which uniqueness witnesses are lifted
through products.

\subsection{Circuits and the structural rules}\label{sec:circuits}

Module \texttt{Circuit.Base} is parameterised by a gate family
$\aid{Gate} : \mathbb{N} \to \mathsf{Set}$:

\begin{agdacode}
data Gen : ℕ → Set where
  gate₀ : Gate 0 → Gen n
  gate₁ : Gate 1 → Gen (₁₊ n)
  gate₂ : Gate 2 → Gen (₂₊ n)
  _↥   : Gen n → Gen (₁₊ n)

Circuit : ℕ → Set
Circuit n = Word (Gen n)
\end{agdacode}

A generator is a gate at the bottom of the wires shifted up by some
\aid{{\_}↥}s; a 0-ary gate is a global scalar, available at every width;
$c{\uparrow} = \aid{wmap}\;\aid{{\_}↥}$ shifts a circuit.  The indices
are built from successors only, so that the coverage checker solves its
unification problems by injectivity of \aid{₁₊} --- which is why the
coset tables below can be defined by direct pattern matching.  Given any
gate-specific family of relations, \aid{Lift-Relation} adds the
structural rules that every circuit theory shares:

\begin{agdacode}
module Lift-Relation (_SRel,_===_ : (n : ℕ) → CRel n) where
  data _VRel,_===_ : (n : ℕ) → CRel n where
    srel  : n SRel, w === v → n VRel, w === v
    cong↑ : n VRel, w === v → (₁₊ n) VRel, w ↑ === v ↑
    comm₀ : (h : Gate 0) (g : Gen n) → n VRel,
      [ g ]ʷ • [ gate₀ h ]ʷ === [ gate₀ h ]ʷ • [ g ]ʷ
    comm₁ : (h : Gate 1) (g : Gen n) → (₁₊ n) VRel,
      [ g ↥ ]ʷ • [ gate₁ h ]ʷ === [ gate₁ h ]ʷ • [ g ↥ ]ʷ
    comm₂ : (h : Gate 2) (g : Gen n) → (₂₊ n) VRel,
      [ g ↥ ↥ ]ʷ • [ gate₂ h ]ʷ === [ gate₂ h ]ʷ • [ g ↥ ↥ ]ʷ
    ω↑=ω : (ω : Gate 0) → (₁₊ n) VRel, [ gate₀ ω ]ʷ ↑ === [ gate₀ ω ]ʷ
\end{agdacode}

Shift-congruence, far-commutation of disjoint gates, and the
identification of a scalar with its shifts: these are the ``monoidal''
rules that the qupit paper leaves implicit in the phrase ``congruence
compatible with composition and tensor product'', and every rule set in
this paper is such a lift.  The lemma \aid{lemma-cong↑} lifts the entire
congruence one wire up, so equational developments move freely along the
wire chain $\mathsf{Cir}\,k \le \mathsf{Cir}\,(k{+}1)$.

\subsection{Coset tables and towers}\label{sec:engine}

The library's normal forms all follow one recipe from computational
group theory~\cite{sims1994computation,holt2005handbook}.  Given a chain
$1 = G_0 \le G_1 \le \dots \le G_n = G$ with a transversal $T_k$ of each
step, every element factors uniquely as $t_1 t_2 \cdots t_n$ --- a
mixed-radix numeral, one digit per level.  For circuits the chain is
handed to us by the wire structure: $\mathsf{Cir}\,k$ embeds into
$\mathsf{Cir}\,(k{+}1)$ as the circuits that do not touch the bottom
wire, and a transversal is a family of \emph{coset representatives}
indexed by $k$.  One level is a Reidemeister--Schreier situation, packaged
as the module \aid{SingleLevel} with parameters
\begin{align*}
  &\Gamma : \aid{WRel}\;X,\quad \Delta : \aid{WRel}\;Y,\quad C : \mathsf{Set},\quad
  I : C,\quad f : X \to \aid{Word}\;Y,\\
  &h : C \to Y \to \aid{Word}\;X \times C,\quad [\_] : C \to \aid{Word}\;Y
\end{align*}
--- a subgroup presentation, a group presentation, the cosets, the
identity coset, the generator embedding, the \emph{coset table} $h$ that
pushes one letter of $G$ past a coset, emitting a word of the subgroup and
the successor coset, and the Schreier section.  Five hypotheses make the
data a genuine coset presentation:

\begin{agdacode}[basicstyle=\footnotesize\ttfamily]
(h=⁻¹f-gen : ∀ (x : X) → ([ x ]ʷ , I) ~ ((h ᵗ) I (f x)))
(h-wd-ax : ∀ (c : C){u t : Word Y} → u ===₂ t → ((h ᵗ) c u) ~ ((h ᵗ) c t))
(f-wd-ax : ∀ {w v} → w ===₁ v → (f ʷ) w ≈₂ (f ʷ) v)
([I]≈ε : [ I ] ≈₂ ε)
(h=ract :  ∀ c b → let (b' , c') = h c b in
  [ c ] • [ b ]ʷ ≈₂ (f ʷ) b' • [ c' ])
\end{agdacode}

From these, \aid{Transfer} derives the normal-form transport
\[
  \aid{nfp'} : \aid{NormalForm}\;\Gamma\;\NF \to
  \aid{NormalForm}\;\Delta\;(\NF \times C) :
\]
the new \aid{nf} runs the
table from the identity coset and normalises the accumulated subgroup
word, the new section is $(f\,\aid{ʷ})(\aid{inv-nf}\;u) \bullet [c]$,
and the hypotheses prove they are inverse.  \aid{CosetTower} iterates
this up an $\mathbb{N}$-indexed family of presentations, one
\aid{Extension} record per level, with carrier
\[
  \aid{tower-carrier}\;\NF_0\;(k{+}1) = \aid{tower-carrier}\;\NF_0\;k \times C_k .
\]
The hypotheses are \emph{quantifier-small} --- they range over cosets,
generators and axioms, all finite inductive types --- and in the symmetric
group they close by evaluation.  In the symplectic group they do not, and
\S\ref{sec:pushing} explains what replaces evaluation there.

\subsection{Presentation construction theorems}\label{sec:constructions}

Relation sets over a disjoint union of alphabets are built from one
primitive, a three-way join $\Gamma_1 \aid{⋄} \Gamma_2 \aid{⋄} \Gamma_3$
whose \emph{mixed} component $\Gamma_3 : \aid{WRel}\;(A \uplus B)$ is
the design parameter.  Two mixed components matter here.  For the
semidirect product the mixed relation says that moving a letter $h$ of
the acting group past a letter $n$ of the normal subgroup conjugates
$n$, by a \emph{word}:

\begin{agdacode}
data ConjRelʷ {N H} (conj : H → N → Word N) : WRel (N ⊎ H) where
  comm : (n : N) (h : H) →
         ConjRelʷ conj ([ [ h ]ʷ ]ᵣ • [ [ n ]ʷ ]ₗ)
                  ([ conj h n ]ₗ • [ [ h ]ʷ ]ᵣ)
\end{agdacode}

The semidirect presentation theorem is parameterised by the two relation
sets and the action, by two well-definedness conditions on the action ---
it respects the acting group's axioms in its first argument and the
normal subgroup's axioms in its second --- and by presentations of the
two factors:

\begin{agdacode}
module Presentation.Construct.Properties.SemiDirectProduct
  {N H : Set} (Γ : WRel N) (Δ : WRel H) (conj : H → N → Word N) where
module _
  (conj-hyph : ∀ {c d} n → c ===₂ d → (conj ʰ') c n ≈₁ (conj ʰ') d n)
  (conj-hypn : ∀ c {w v} → w ===₁ v → (conj ⁿ') c w ≈₁ (conj ⁿ') c v)
  where
  module Presentation (G1 G2 : Group 0ℓ 0ℓ)
    (p1 : Γ IsPresentationOf G1) (p2 : Δ IsPresentationOf G2) where
    dpres : (Γ ⋄ Δ ⋄ ConjRelʷ conj) IsPresentationOf G1⋊G2
\end{agdacode}

Here \aid{G1⋊G2} is the semidirect product of the two \emph{semantic}
groups, built by the library's \texttt{ForStdlib} construction from a
bundled \aid{Action} record whose action is the syntactic \aid{conj}
transported through the two presentation isomorphisms.  Its normal form
is the product of the factors' normal forms; the proof is a
Reidemeister--Schreier level whose cosets are words over the acting
alphabet up to $\approx$, and its uniqueness is lifted from the factors
by \aid{by-completeness}.  The second mixed component,
\aid{RelTwist}, twists each relator $u = v$ of the quotient by a
correction word in the normal subgroup, $[u]_r = [\mathit{corr}\;r]_l
\bullet [v]_r$, and
\begin{agdacode}
extension-presentation S R conj corr = S ⋄ EmptyRel ⋄ (ConjRelʷ conj ∪ RelTwist R corr)
\end{agdacode}
is the recipe of Proposition~2.55 of the Handbook of Computational
Group Theory~\cite{holt2005handbook}, which the qupit paper invokes for
its Section~4.2: the relations of the normal subgroup, conjugation rules,
and each relator of the quotient corrected by the element of the normal
subgroup by which its lift fails to hold.  The library proves it as a
presentation theorem for any group extension (given a short exact
sequence realising the data) and, in the form used in
\S\ref{sec:scalars}, for central extensions built from a cocycle.  The
semidirect product is the special case with trivial corrections.

\subsection{Modular arithmetic and primitive roots}\label{sec:modp}

Finally, $\Zp$ itself.  \texttt{ForStdlib.Data.Fin.Mod} models $\Zp$ as
$\mathsf{Fin}\;p$ with operations on representatives, $\Zp^{*}$ as the
nonzero residues, inversion by B\'ezout, and a ring solver; the module
\aid{PrimeModulus'} adds Fermat's little theorem, proved by the classical
argument that multiplication by a unit permutes the nonzero residues,
and \aid{Primitive-Root-Modp'} takes a primitive root $g$ and the
witness \aid{g-gen} as parameters and derives from them the discrete
logarithm and the exponent-reduction law $g^{k} = g^{k \bmod (p-1)}$.
These are the arithmetic facts behind the one-qupit rules: the
multiplier law $M_x M_y = M_{xy}$ is, on exponents, addition modulo
$p-1$, and Fermat is what makes $M_g^{\,p-1} = \varepsilon$.
